\chapter{Chest Pain}

\section*{Definition and Overview}
Chest pain is a common symptom that can arise from the heart, lungs, oesophagus, muscles, or ribs, and it may also be felt as referred pain from structures outside the chest. Its significance ranges from minor muscle strain to life-threatening emergencies such as a heart attack, a blood clot in the lung, or a tearing of the aorta. The central task in assessing chest pain is to distinguish urgent cardiac and lung causes from the many less serious conditions that produce similar discomfort. Anyone with concerning features should be assessed promptly, because delays in treating a heart attack or other emergency can be fatal.

\section*{Epidemiology}
Chest pain is among the most frequent reasons for emergency department attendance in adults. Coronary heart disease remains a leading cause of death, and chest pain is often its first warning. Older people, men, smokers, and people with diabetes, high blood pressure, or high cholesterol are at greater risk of a cardiac cause. Younger people and those without cardiovascular risk factors more often have musculoskeletal, gastrointestinal, or anxiety-related chest pain. Nevertheless, serious causes such as pulmonary embolism and aortic dissection can occur at any age, which is why chest pain is always taken seriously.

\section*{Aetiology and Pathophysiology}
\begin{itemize}
    \item \textbf{Cardiac:} angina and heart attack occur when narrowed or blocked coronary arteries reduce blood flow to the heart muscle
    \item \textbf{Pericarditis:} inflammation of the sac surrounding the heart
    \item \textbf{Aortic dissection:} a tear in the wall of the aorta, the large artery leaving the heart
    \item \textbf{Pulmonary:} pulmonary embolism, pneumonia, pleurisy, or a collapsed lung (pneumothorax)
    \item \textbf{Gastrointestinal:} acid reflux, oesophagitis, or spasm of the oesophagus
    \item \textbf{Musculoskeletal:} costochondritis, rib injury, or strain of the chest wall muscles
    \item \textbf{Psychological:} anxiety and panic attacks producing chest tightness and pain
\end{itemize}

\section*{Clinical Features}
\begin{itemize}
    \item \textbf{Angina and heart attack:} crushing or heavy central chest pain, often spreading to the arm, neck, or jaw, with sweating and breathlessness
    \item \textbf{Pulmonary embolism:} sudden sharp chest pain with shortness of breath, and sometimes coughing of blood
    \item \textbf{Pericarditis:} sharp pain worse when lying flat or breathing in, often eased by sitting forward
    \item \textbf{Reflux:} a burning discomfort behind the breastbone, often after meals or when lying down
    \item \textbf{Musculoskeletal:} localised tenderness reproduced by pressing on the chest wall or by certain movements
    \item \textbf{Panic attack:} chest tightness with rapid breathing, dizziness, and a sense of dread
    \item \textbf{Aortic dissection:} sudden, severe tearing or ripping pain, often radiating to the back
\end{itemize}

\section*{Differential Diagnosis}
\begin{itemize}
    \item Heart attack and unstable angina
    \item Aortic dissection
    \item Pulmonary embolism and pneumothorax
    \item Pericarditis and myocarditis
    \item Pneumonia and pleurisy
    \item Gastro-oesophageal reflux disease and oesophageal spasm
    \item Costochondritis, rib fracture, and intercostal muscle strain
    \item Panic disorder and anxiety
    \item Shingles (herpes zoster) before the rash appears
\end{itemize}

\section*{When to Seek Medical Care}
Seek same-day medical advice for new or worsening chest pain, unexplained discomfort, or pain that occurs with exercise and settles with rest, since these can be features of angina.

\medskip
\textbf{Call 000 (or local emergency services) for:}
\begin{itemize}
    \item Sudden, severe, crushing, or squeezing chest pain lasting more than a few minutes
    \item Chest pain that spreads to the arm, back, neck, or jaw
    \item Chest pain with sweating, nausea, shortness of breath, or feeling faint
    \item Sudden onset of breathlessness or coughing up blood
    \item Tearing or ripping chest pain, especially radiating to the back
\end{itemize}

\section*{Diagnosis and Investigations}
\begin{itemize}
    \item \textbf{History:} the character, site, onset, duration, and triggers of the pain, and what relieves it
    \item \textbf{Physical examination:} pulse, blood pressure in both arms, and examination of the heart and lung sounds
    \item \textbf{Electrocardiogram (ECG):} a tracing of the heart's electrical activity to detect a heart attack or rhythm disturbance
    \item \textbf{Blood tests:} cardiac enzymes such as troponin, a full blood count, and markers of inflammation
    \item \textbf{Chest X-ray:} to look for lung, heart, or aortic problems
    \item \textbf{Stress testing:} an ECG while exercising to uncover reduced blood flow to the heart
    \item \textbf{CT angiography:} detailed scans to detect pulmonary embolism or aortic dissection
    \item \textbf{Echocardiogram:} ultrasound of the heart to assess function and look for pericardial fluid
\end{itemize}

\section*{Management}
\begin{itemize}
    \item \textbf{Emergency heart attack treatment:} aspirin, nitrates, oxygen, and urgent hospital care with reperfusion therapy
    \item \textbf{Angina:} medicines to relieve symptoms and reduce risk, together with rigorous control of risk factors
    \item \textbf{Anticoagulants:} for pulmonary embolism and other clot-related causes
    \item \textbf{Acid-reducing medicines:} for reflux and oesophagitis, alongside dietary changes
    \item \textbf{Anti-inflammatory medicines:} for pericarditis and musculoskeletal causes
    \item \textbf{Relaxation and psychological support:} for anxiety-related chest pain
    \item \textbf{Follow-up:} repeat assessment to confirm the diagnosis and adjust treatment
\end{itemize}

\section*{Complications}
\begin{itemize}
    \item Missed heart attack, with subsequent damage to heart muscle
    \item Cardiac arrest or life-threatening rhythm disturbance
    \item Pulmonary embolism causing respiratory or circulatory failure
    \item Aortic dissection leading to severe internal bleeding
    \item Chronic pain and reduced quality of life when the cause remains unexplained
    \item Persistent anxiety about the heart even after investigations are normal
\end{itemize}

\section*{Prognosis and Outcomes}
The outlook depends entirely on the underlying cause. Muscle and rib causes usually settle within days to weeks. Angina can be well controlled with medication, but a heart attack can be life-threatening if treatment is delayed, which is why rapid assessment matters. With modern treatment, most people survive heart attacks and return to a good quality of life. Recurrent chest pain of any cause should always be reviewed by a doctor, since symptoms and risk can change over time.

\section*{Prevention and Self-Care}
\begin{itemize}
    \item Do not smoke, and avoid second-hand smoke
    \item Manage blood pressure, cholesterol, and diabetes with medical advice
    \item Eat a heart-healthy diet low in salt and saturated fat
    \item Aim for regular physical activity as advised by a doctor
    \item Maintain a healthy weight and limit alcohol
    \item Learn and practise the warning signs of a heart attack
    \item Take prescribed cardiac medicines exactly as directed
\end{itemize}

\section*{Sources and Further Reading}
The following reputable sources provide further information on chest pain:
\begin{itemize}
    \item NHS. \textit{Health A--Z: chest pain information.} https://www.nhs.uk/conditions/
    \item Mayo Clinic. \textit{Diseases and conditions library.} https://www.mayoclinic.org/
    \item British Heart Foundation. \textit{Heart attack and chest pain information.} https://www.bhf.org.uk/
    \item NICE. \textit{Clinical guidelines on chest pain and acute coronary syndromes.} https://www.nice.org.uk/
    \item MSD Manual. \textit{Professional edition: chest pain evaluation.} https://www.msdmanuals.com/
    \item MedlinePlus. \textit{Health topics on chest pain and heart conditions.} https://medlineplus.gov/
\end{itemize}
